\documentclass[12pt,a4paper]{article}
\usepackage[UTF8]{ctex}
\usepackage[margin=2cm,top=2cm,bottom=2cm,left=2.5cm,right=2.5cm]{geometry}
\usepackage{setspace}
\usepackage{indentfirst}
\usepackage{cite}
\usepackage{amsmath}
\usepackage{amssymb}

\setlength{\parindent}{2em}
\onehalfspacing

\begin{document}

\title{\textbf{网络零售商直播模式选择策略的文献综述}}
\author{}
\date{}
\maketitle

\section{引言}

近年来，直播电商作为一种新兴的数字营销模式迅速崛起，成为网络零售商拓展销售渠道、提升品牌影响力的重要手段。据商务部数据显示，2024年全国直播电商市场规模已突破万亿元，直播带货已成为电商行业的核心增长引擎。然而，网络零售商在引入直播渠道时面临着关键的策略选择问题：是采用商家自播模式（由企业员工担任主播），还是采用网红直播模式（聘请网络知名人士担任主播）？两种模式在成本结构、销售努力程度、消费者互动效果等方面存在显著差异，如何科学地选择直播模式已成为学术界和实务界共同关注的焦点。

与此同时，在线评论作为消费者购买决策的重要信息来源，对直播模式选择策略产生着深远影响。负面在线评论可能抑制消费者购买意愿，而主播可通过针对性话术缓解负面评论的不利影响，但不同直播模式下主播的销售努力程度及其对负面评论的缓解效果各不相同。因此，在线评论信息与主播佣金率的交互作用如何影响网络零售商的直播模式选择，是一个值得深入探讨的理论与现实问题。

本文旨在对网络零售商直播模式选择策略的相关研究进行系统梳理，从直播电商与消费者行为、在线评论与定价渠道决策、直播渠道引入策略、直播模式选择策略四个维度展开文献回顾，评述现有研究的理论贡献与不足，并展望未来研究方向。

\section{直播电商与消费者行为}

直播电商的本质是通过实时视频互动缩短消费者与产品之间的信息距离，从而提升购买转化率。彭宇泓等（2021）研究了直播营销中关系纽带与顾客承诺对消费者在线购买意愿的影响，发现直播中的互动关系能够显著增强消费者的信任感和购买意愿\cite{peng2021}。PIG等（2023）从最优销售格式视角考察了直播销售中的价格歧视策略，揭示了直播渠道中差异化定价的可行性与效率\cite{pig2023}。HUANG等（2024）则关注传统零售商是否应引入直播渠道的问题，通过渠道竞争模型分析了两类零售商的均衡策略\cite{huang2024}。WANG等（2023）研究了直播销售的战略引入问题，发现直播渠道的引入有助于扩大市场覆盖面，但需权衡渠道间的蚕食效应\cite{wang2023}。

在消费者行为层面，刘凤军等（2021）从双面信息视角研究了在线产品评论负面评语的"明亮面"，发现适度的负面评论反而能增强消费者对评论真实性的感知，从而提升购买意愿\cite{liu2021}。FAN等（2022）研究了在线评论和直播销售对消费者购买决策的交互影响，发现主播可从在线评论中挖掘消费者偏好信息，并通过针对性话术缓解负面评论的不利影响\cite{fan2022}。这些研究为理解直播电商情境下消费者行为的变化规律奠定了理论基础。

\section{在线评论对定价与渠道结构决策的影响}

在线评论作为消费者生成内容（UGC）的重要形式，对零售商的定价决策和渠道结构决策产生着广泛影响。该研究方向可进一步细分为在线评论对定价策略的影响和对渠道结构决策的影响两个方面。

\subsection{在线评论对定价策略的影响}

在线评论所揭示的产品质量信息、平均评级、评论数量等维度均会影响零售商的定价决策。WANG等（2022）研究了供应商在存在消费者评论时的定价策略，发现评论的情感倾向显著影响最优定价方向\cite{wang2022}。李昊等（2026）基于下游竞争环境研究了在线评论的引入策略，揭示了在线评论对供应链定价的溢出效应\cite{li2026}。TANG等（2025）将在线评论纳入先进销售策略与定价决策框架，发现评论的平均星级对预售定价具有显著影响\cite{tang2025}。ZHAO等（2024）研究了在线评论数量对双渠道定价的影响，发现评论数量的增加会加剧渠道间的价格竞争\cite{zhao2024}。

\subsection{在线评论对渠道结构决策的影响}

在线评论的积极程度与有效程度也会影响零售商的渠道结构选择。郑本荣等（2022）构建了考虑在线评论的双渠道供应链定价与服务决策模型，发现在线评论的积极程度正向影响零售商引入线上渠道的动机\cite{zheng2022}。黄鹤（2023）研究了在线评论对电商平台定价与消费者渠道选择的影响，揭示了评论信息在渠道迁移中的桥梁作用\cite{huang2023}。金亮和郑本荣（2023）从在线评论价值视角研究了电商平台自营渠道引入决策，发现评论带来的额外价值是驱动平台引入自营渠道的关键因素\cite{jin2023}。YE等（2022a）考察了在线评论对零售商渠道策略的影响，提出了"点击与 mortar"渠道与线上渠道的最优组合条件\cite{ye2022a}。YE等（2022b）进一步研究了在线评论对传统零售商与线上零售商渠道结构决策的影响，发现评论信息的溢出效应改变了零售商的渠道选择均衡\cite{ye2022b}。张艳芬等（2023）从产品评论价值视角研究了供应链激励契约设计，为在线评论的价值量化提供了分析框架\cite{zhang2023}。

\section{直播渠道引入策略}

直播渠道引入策略的研究关注零售商在何种条件下应引入直播渠道，以及引入直播渠道后对传统渠道的影响。该方向的研究从早期关注直播渠道的单纯引入决策，逐步扩展到考虑多种影响因素的综合性决策。

ZHANG等（2026）通过构建直播销售渠道选择的分析框架，比较了零售商在不同渠道结构下的均衡利润，为直播渠道引入决策提供了博弈论基础\cite{zhang2026}。王玉燕等（2025）研究了直播电商供应链中制造商的直播引入策略，发现主播佣金率是决定制造商是否引入直播渠道的关键因素\cite{wang2025}。WANG等（2025）研究了第三方直播电商模式下的渠道竞争问题，揭示了直播渠道引入对传统渠道的蚕食效应与市场扩张效应的权衡关系\cite{wang2025b}。

LU等（2025）从零售商信息共享与制造商渠道选择视角研究了直播渠道引入决策，发现信息共享程度和直播渠道的溢出效应共同决定了渠道引入的均衡结果\cite{lu2025}。PAN等（2025）研究了直播平台的产品推荐策略，发现推荐算法的精准度直接影响零售商引入直播渠道的意愿\cite{pan2025}。DU等（2023）从最优直播渠道选择的物流视角分析了运输模式的选择问题，将物流成本纳入直播渠道引入的决策框架\cite{du2023}。LI等（2025）研究了零售商补贴的附加优惠策略对直播渠道引入的影响，发现消费者对直播渠道的接受程度是引入决策的关键约束\cite{li2025}。汪乐等（2024）研究了制造商直播引入策略，将消费者对产品的感知价值纳入模型，发现感知价值较低时引入直播渠道的动机更强\cite{wang2024}。

\section{直播模式选择策略}

在确认引入直播渠道后，零售商还需在商家自播与网红直播两种模式之间进行选择。这是直播电商研究中最具策略性的问题之一，涉及主播佣金、销售努力、消费者渠道偏好等多重因素的权衡。

ZHANG和TANG（2023）从制造商是否应开设直播渠道的视角研究了直播模式选择，发现当网红主播佣金率较低时制造商倾向于选择网红直播，反之则倾向于自播\cite{zhang2023}。张鑫和张杰（2025）研究了直播背景下电商供应链混合渠道定价与选择策略，揭示了不同直播模式下最优定价策略的差异\cite{zhangxin2025}。YANG等（2023）从直播销售的双渠道运营视角研究了供应链均衡，发现商家自播模式下的供应链协调效果优于网红直播模式\cite{yang2023}。FAN等（2024）进一步考察了直播渠道中互动效应对直播模式选择的影响，发现互动性越强，网红直播模式的相对优势越明显\cite{fan2024}。LI等（2024）研究了直播渠道选择的消费者行为模型，发现消费者的渠道忠诚度显著影响直播模式的选择偏好\cite{li2024}。陈悦等（2026）将制造商入侵纳入直播销售策略研究，发现当存在入侵威胁时零售商更倾向于选择商家自播以保持渠道控制权\cite{chen2026}。

尤天慧等（2026）针对考虑在线评论的网络零售商直播模式选择问题，通过构建不引入直播模式、商家自播模式和网红直播模式下的Stackelberg博弈模型，综合考察了在线评论信息和主播佣金率对直播模式选择策略的交互影响\cite{you2026}。研究发现：随着在线评论信息趋于负面，零售商更倾向于引入直播渠道；若引入直播，选择何种直播模式受在线评论信息和主播佣金率的共同影响。该研究将在线评论与直播渠道决策纳入统一分析框架，弥补了既有研究鲜少涉及在线评论与直播销售交互的不足，为网络零售商的直播模式选择提供了重要的理论依据和管理启示。

\section{评述与展望}

综合以上文献回顾，现有研究在网络零售商直播模式选择策略方面已取得较为丰富的成果，但仍存在以下不足和值得深入探索的方向：

第一，现有研究大多将在线评论视为外生变量，未充分考虑在线评论与直播销售的动态交互机制。实际上，直播过程中主播的话术和互动可能引发消费者产生新的在线评论，形成"直播$\rightarrow$评论$\rightarrow$购买决策"的动态反馈循环。未来研究可构建动态博弈模型，考察在线评论的内生演化及其对直播模式选择的长期影响。

第二，现有文献在刻画直播模式差异时，主要从佣金率、销售努力等供给侧维度进行区分，对需求侧因素（如消费者对不同直播模式的心理感知差异、信任异质性等）关注不足。未来研究可将消费者的直播模式偏好、信任差异等行为因素纳入模型，更精细地刻画两种直播模式下的消费者响应机制。

第三，多数研究假设市场结构为垄断或双边垄断，对竞争性市场环境下的直播模式选择策略研究较为缺乏。现实中，多个网络零售商可能同时引入直播渠道，形成直播渠道间的横向竞争。未来研究可拓展至多零售商竞争框架，考察竞争强度对直播模式选择均衡的影响。

第四，现有研究对直播平台的平台策略（如流量分配、佣金规则制定等）关注较少。作为连接零售商与消费者的双边平台，直播平台的策略选择直接影响零售商的直播模式决策。未来研究可从平台经济学视角，考察平台定价、流量倾斜等策略对零售商直播模式选择的间接影响，丰富平台治理的理论内涵。

第五，从方法论角度看，现有研究以博弈论建模为主，实证研究相对匮乏。未来可利用直播电商平台的微观数据，采用计量经济学方法实证检验在线评论、主播特征等因素对直播模式选择和销售绩效的影响，弥补理论模型与实证证据之间的缺口。

\begin{thebibliography}{99}

\bibitem{peng2021}
彭宇泓, 韩欢, 郝辽钢, 等. 直播营销中关系纽带、顾客承诺对消费者在线购买意愿的影响研究[J]. 管理学报, 2021, 18(11): 1686-1694.

\bibitem{pig2023}
PIG P, FU T Y, LI S N. Optimal selling format considering price discrimination strategy in live-streaming[J]. European Journal of Operational Research, 2023, 309(2): 529-544.

\bibitem{huang2024}
HUANG L C, LIU B, ZHANG R. Channel strategy for competing retailers: whether and when to introduce live-streaming?[J]. European Journal of Operational Research, 2024, 312(2): 413-426.

\bibitem{wang2023}
WANG S S, GUO X. Strategic introduction of live-streaming in a funneling platform[J]. Electronic Commerce Research and Applications, 2023, 62: 101315.

\bibitem{liu2021}
刘凤军, 段鑫, 孟陆, 等. 瑕不掩瑜?在线产品评论负面评语的明亮面——基于双面信息视角研究[J]. 管理工程学报, 2021, 35(5): 89-101.

\bibitem{fan2022}
FAN X P, TAO W X, YU N W, et al. Considering online reviews and live-streaming sales' impact on consumer purchase decision: an online/O2O channel strategy[J]. Production and Operations Management, 2022, 31(9): 3387-3399.

\bibitem{wang2022}
WANG P Y, SHU B S, FENG G Z. Supplier's pricing strategy in the presence of consumer reviews[J]. European Journal of Operational Research, 2022, 296(2): 570-586.

\bibitem{li2026}
李昊, 晏妮娜, 卢继周, 等. 基于下游竞争的在线评论引入策略研究[J]. 中国管理科学, 2026, 34(3): 104-113.

\bibitem{tang2025}
TANG M Z, YOU T N, CAO B B. Advance selling strategy and pricing decisions with online reviews[J]. International Transactions in Operational Research, 2025, 32(5): 3108-3137.

\bibitem{zhao2024}
ZHAO N N, XU P Y. Pricing strategy considering online reviews in a dual-channel supply chain[J]. Managerial and Decision Economics, 2024, 45(5): 3254-3267.

\bibitem{zheng2022}
郑本荣, 李芯怡, 黄燕婷. 考虑在线评论的双渠道供应链定价与服务决策[J]. 管理学报, 2022, 19(2): 289-298.

\bibitem{huang2023}
黄鹤. 在线评论对电商平台定价与消费者渠道选择的影响[J]. 管理评论, 2023, 35(12): 148-159.

\bibitem{jin2023}
金亮, 郑本荣. 电商平台自营渠道引入决策: 在线评论产品评论的价值[J]. 系统工程理论与实践, 2023, 43(2): 469-487.

\bibitem{ye2022a}
YE F, LIANG L N, TONG Y, et al. Brick-and-mortar or click? The influence of online reviews on retailer's channel strategy[J]. IISE Transactions, 2022, 54(12): 1199-1210.

\bibitem{ye2022b}
YE F, LIANG L N, TONG Y. Impact of online reviews on brick-and-mortar retailers' online channel expansion[J]. International Journal of Production Economics, 2022, 248: 108534.

\bibitem{zhang2023}
张艳芬, 徐琪, 孙中苗. 考虑主播带货努力与影响力的直播电商供应链激励契约研究[J]. 管理学报, 2023, 20(2): 278-286.

\bibitem{zhang2026}
ZHANG W, YU L L, WANG Z Z. Live-streaming sales channel choice: an analysis of strategic manufacturer-retailer interaction[J]. European Journal of Operational Research, 2026, 359: 1745-1769.

\bibitem{wang2025}
王玉燕, 丁露萍, 霍宝锋. 直播电商供应链中制造商的直播引入策略研究[J]. 系统科学与数学, 2025, 45(5): 1471-1493.

\bibitem{wang2025b}
WANG S S, WANG P B. Third-party live-streaming e-commerce model selection under channel competition: a third-party or in-house streamer?[J]. Electronic Commerce Research, 2025, 25(4): 3261-3282.

\bibitem{lu2025}
LU W P, PI X W U P. Retailer information sharing and manufacturer's channel selection in the era of live-streaming e-commerce[J]. European Journal of Operational Research, 2025, 320(3): 527-543.

\bibitem{pan2025}
PAN R, FENG P, ZHAO Z L. Fulfilling watching live-streamer product recommendation strategy[J]. Production and Operations Management, 2025: 1-19.

\bibitem{du2023}
DU Z, FAN Z P, SUN M F N, et al. Optimal streaming sales channel strategy: an analysis of strategic manufacturer-retailer interaction[J]. Transportation Research Part E: Logistics and Transportation Review, 2023, 173: 103096.

\bibitem{li2025}
LI Q Y, ZHU S, CHOI T M, et al. Managing commerce fulfilment sustainability: an agile strategic live-streaming channel[J]. Omega, 2025, 133: 103276.

\bibitem{wang2024}
汪乐, 宋杨, 范体军. 制造商直播引入的策略研究[J]. 中国管理科学, 2024, 32(2): 276-284.

\bibitem{zhang2023b}
ZHANG T, TANG Z P. Should manufacturers open a live-streaming sales channel?[J]. Journal of Retailing and Consumer Services, 2023, 71: 103229.

\bibitem{zhangxin2025}
张鑫, 张杰. 直播背景下电商供应链混合渠道定价与选择策略[J]. 系统管理学报, 2025, 34(1): 27-39.

\bibitem{yang2023}
YANG W T, GOVINDAN K, ZHANG J T. Supply chain equilibrium of live-streaming selling dual-channel operations[J]. Transportation Research Part E: Logistics and Transportation Review, 2023, 180: 103298.

\bibitem{fan2024}
FAN X P, ZHANG L, GUO X, et al. The impact of live-streaming interactive activity on sales[J]. Journal of the Operational Research Society, 2024, 81: 103981.

\bibitem{li2024}
LI Y N, NING Y, FAN W G, et al. Channel choice of live-streaming commerce[J]. Production and Operations Management, 2024, 33(11): 2221-2240.

\bibitem{chen2026}
陈悦, 王勇, 郑静, 等. 考虑制造商入侵的中小型零售商直播销售策略价值研究[J]. 管理学报, 2026, 23(1): 157-167.

\bibitem{you2026}
尤天慧, 李富禄, 曹兵兵. 考虑在线评论的网络零售商直播模式选择策略[J]. 管理学报, 2026.

\end{thebibliography}

\end{document}
